\begin{proposition}
\label{proposition-fppf-open}
Let $R \to S$ be flat and of finite presentation.
Then $\Spec(S) \to \Spec(R)$ is open.
More generally this holds for any ring map $R \to S$ of
finite presentation which satisfies going down.
\end{proposition}

\begin{proof}
If the original map $\varphi:R\to S$ is flat, it
satisfies going down by Lemma \ref{lemma-flat-going-down}.
Thus to prove the proposition it suffices to retain
finite presentation and going down.
For an original $g\in S$, write a finite presentation
$S=R[x_1,\ldots,x_n]/(h_1,\ldots,h_m)$ and choose
an original polynomial $G(x_1,\ldots,x_n)$ representing
$g$. Then the localization has the finite presentation
$$
S_g\simeq
R[x_1,\ldots,x_n,t]/(h_1,\ldots,h_m,tG-1).
$$
The map sends $x_i$ to its original class and $t$
to $1/g$. The inverse sends $s/g^d$ to the class
of $P(x)t^d$ for a polynomial lift $P$ of $s$.
Independence follows from the original relations
$h_j$ and the relation $tG=1$; these maps are
inverse on the algebra generators and fractions.
Hence Chevalley's Theorem \ref{theorem-chevalley}
applies to this particular $R\to S_g$ and shows
that the image $E_g$ of $D(g)$ is constructible.

Going down shows that $E_g$ is stable under
generalization. Given $\mathfrak q'\in D(g)$ over
$\mathfrak p'$ and $\mathfrak p\subseteq\mathfrak p'$
in $R$, take $\mathfrak q\subseteq\mathfrak q'$
over $\mathfrak p$. Since $g\notin\mathfrak q'$,
also $g\notin\mathfrak q$, and the lift remains
in the original $D(g)$. This is also the exact
going-down comparison for $S\to S_g$ and the
composition lemma. The preceding constructible-set
lemma now makes $E_g$ open. Every original open
of $\Spec(S)$ is a union of such $D(g)$ and images
commute with unions, proving openness of the map.

The proof only needs constructibility of the
images of all principal opens. Under that weaker
assumption, going down is equivalent to openness:
the proof just given gives one direction, and
Lemma \ref{lemma-open-going-down} gives the other.
In particular this is an equivalence for maps of
finite presentation. Finite presentation is not
necessary: for a field $k$, the map
$k\to k[x_i:i\geq1]$ has open spectrum map,
because the target spectrum $\Spec(k)$ has one
point. The algebra is not even of finite type:
finitely many polynomials use finitely many
variables and cannot generate an omitted original
indeterminate, as evaluation fixing those variables
and changing that indeterminate verifies.

For the flat finite-presentation case the map is
universally open. Indeed for every original
$R$-algebra $A$, the base change has presentation
$$
S\otimes_RA\simeq
A[x_1,\ldots,x_n]/(h_1^A,\ldots,h_m^A),
$$
where $h_j^A$ retains every coefficient through
the original map $R\to A$. The map sends
$s\otimes a$ to $a$ times a polynomial representative
of $s$; the inverse sends the original variables
to $x_i\otimes1$ and each $a$ to $1\otimes a$.
The displayed relations and tensor balancing make
these maps inverse. Flatness is preserved because
for an injection of $A$-modules $L\hookrightarrow L'$,
the original comparison
$$
L\otimes_A(S\otimes_RA)\simeq L\otimes_RS,
\qquad l\otimes(s\otimes a)\longmapsto al\otimes s,
$$
with inverse $l\otimes s\mapsto l\otimes(s\otimes1)$,
identifies its tensor map with an injection by
flatness of $S$ over $R$. Thus the proved proposition
applies after every base change. This last assertion
uses flatness, not merely the original going-down
hypothesis.
\end{proof}